\section{La esencia del Agente: el ciclo ReAct, las llamadas a herramientas y el cómputo en tiempo de prueba}

\subsection{De buzzword a la definición mínima}
Sutton, hacia el minuto 00:04, da directamente su propia definición: lo que llamamos un Agente LLM es, en esencia, \textbf{cómputo dinámico + uso de herramientas}. El cómputo dinámico significa que el modelo no produce la respuesta final de una sola vez, sino que asigna más cómputo en tiempo de prueba según la dificultad de la tarea; el uso de herramientas significa que puede externalizar sus hipótesis intermedias hacia un entorno ejecutable y corregir continuamente su siguiente estrategia mediante comandos, ejecución de código e información de depuración.

\begin{knowledgebox}{El ciclo mínimo de un Agente ejecutable}
\[
\text{Trajectory}_{t}
\xrightarrow{\text{LLM}}
\text{Thought + Tool Call}
\xrightarrow{\text{Environment}}
\text{Observation}_{t+1}
\xrightarrow{\text{append}}
\text{Trajectory}_{t+1}
\]
Este ciclo coincide con la idea de ReAct: el razonamiento no es el punto final, sino lo que permite decidir la siguiente acción con una herramienta.
\end{knowledgebox}

\begin{importantbox}{Por qué el «cómputo dinámico» es más determinante que la «inteligencia de una sola vez»}
\begin{itemize}
    \item Las tareas complejas suelen cometer errores desde los primeros pasos; solo permitir la iteración expone y corrige las rutas equivocadas.
    \item El mismo modelo, bajo ``múltiples rondas de prueba y error + retroalimentación ejecutable'', tiene un techo de capacidad notablemente más alto que bajo ``una sola generación''.
    \item Para problemas de ingeniería reales, la verificación continua vale más que la fluidez de un solo texto.
\end{itemize}
\end{importantbox}

Esta definición coincide con la frase que se repite en la clase: \textit{``Chain of Thought plus tool use is at the heart of agents.''} En escenarios de coding, la salida de la herramienta es la señal de entrenamiento para la siguiente ronda de razonamiento; en escenarios de seguridad, el depurador, el sanitizer y los registros de fallos del programa se convierten en una retroalimentación externa más fuerte.

\begin{figure}[H]
\centering
\includegraphics[width=0.88\textwidth]{frame-07-react-loop.jpg}
\caption{En la parte final del curso se retoma el ciclo de estilo ReAct: observar, analizar, llamar a la herramienta y observar de nuevo \protect\footnotemark}
\end{figure}
\footnotetext{Intervalo de tiempo en el video: 01:18:06--01:18:40.}

\subsection{Presupuesto de cómputo en tiempo de prueba: seguir corrigiendo o reiniciar la trayectoria}
Una pregunta muy práctica en la clase es: cuando se obtiene un patch ``parcialmente correcto pero que no pasa'', ¿conviene seguir corrigiéndolo en la trayectoria actual o descartarlo y volver a muestrear? En realidad esto es \textbf{asignación de presupuesto de cómputo en tiempo de prueba}. Para cambios pequeños, volver a muestrear puede ser más barato; para modificaciones complejas de varios archivos, conservar el contexto y seguir corrigiendo puede ser mejor.

\begin{warningbox}{Error común: mirar solo la calidad de una trayectoria}
Si solo se analiza la mejor trayectoria, es fácil sobrestimar el sistema. En un despliegue real conviene observar al mismo tiempo:
\begin{itemize}
    \item el número promedio de trayectorias que consume cada éxito;
    \item si las trayectorias fallidas son recuperables;
    \item la proporción de fallos por tiempo agotado o por bucles.
\end{itemize}
\end{warningbox}

\subsection{La herramienta no es una función accesoria, sino el límite de la capacidad}
Como se plantea en la clase, la llamada a herramientas no es ``darle al modelo un poco de plugin'', sino que determina directamente el espacio de problemas resolubles. Sin una herramienta de exploración de código, el modelo casi no puede ubicar de manera estable los problemas de un repositorio grande; sin una herramienta de retroalimentación de ejecución, el modelo se queda por mucho tiempo en conjeturas semánticas; sin capacidad de depuración, en escenarios de seguridad la cadena de activación de vulnerabilidades es difícil de cerrar.

\begin{importantbox}{El significado sistémico de la llamada a herramientas}
El techo de capacidad de un Agente no lo determina solo el base model, sino que también lo define en conjunto ``el conjunto de herramientas $\times$ el formato de retroalimentación $\times$ la estrategia de control''. Quien imparte la clase llama a estos tres elementos un diseño a nivel de sistema que debe ajustarse en conjunto.
\end{importantbox}

\subsection{Resumen de la sección}
La definición de Agente en esta clase es clara: es un sistema de acción que puede iterar de manera sostenida, no un generador de respuestas largas. El cómputo dinámico, la retroalimentación de herramientas y las trayectorias recuperables constituyen la base de toda la discusión de ingeniería que sigue.

